\section*{Chapitre \ref{ch:g6:light-propagation} --- La propagation rectiligne de la lumière}

\begin{solution}{exo:g6:light-propagation:1}
Dans l'\omterm{def:g2:air-around-us:air}{air} limpide (ou tout milieu limpide et homogène), la \omterm{def:g1:light-and-shadows:source}{lumière}
se propage en ligne droite. Rectiligne~: en ligne droite~;
propagation~: le voyage de la \omterm{def:g1:light-and-shadows:source}{lumière} d'un endroit à un autre.
\end{solution}

\begin{solution}{exo:g6:light-propagation:2}
Divergent~: des rayons qui s'écartent d'un point (une bougie, une
lampe nue). Parallèle~: des rayons côte à côte (la \omterm{def:g1:light-and-shadows:source}{lumière} du
Soleil --- le Soleil est si loin que ses rayons nous parviennent au
pas). Convergent~: des rayons qui se referment vers un point (de la
\omterm{def:g1:light-and-shadows:source}{lumière} rassemblée par un instrument grossissant).
\end{solution}

\begin{solution}{exo:g6:light-propagation:3}
C'est un rayon de \omterm{def:g1:light-and-shadows:source}{lumière} qui joue la règle~: l'œil, l'arête et
l'extrémité lointaine sont jugés alignés exactement quand la \omterm{def:g1:light-and-shadows:source}{lumière}
venue de l'extrémité atteint l'œil en frôlant toute l'arête --- la
\omterm{prop:g6:light-propagation:law}{propagation rectiligne} est le serment.
\end{solution}

\begin{solution}{exo:g6:light-propagation:4}
Les points d'arrivée des rayons qui frôlent tout juste les bords de
l'obstacle. Un œil placé à l'intérieur du cône d'\omterm{def:g1:light-and-shadows:shadow}{ombre} ne peut pas
voir la lampe --- aucun rayon rectiligne issu de la lampe ne
l'atteint.
\end{solution}

\begin{solution}{exo:g6:light-propagation:5}
À mi-chemin~: l'\omterm{def:g1:light-and-shadows:shadow}{ombre} fait deux fois la balle --- \qty{12}{cm}. Près
de l'écran~: l'\omterm{def:g1:light-and-shadows:shadow}{ombre} se resserre vers les \qty{6}{cm} réels de la
balle.
\end{solution}

\begin{solution}{exo:g6:light-propagation:6}
En bas de l'écran~: une ligne droite issue d'un point haut, passant
par la porte plus basse, doit continuer vers le bas.
\end{solution}

\begin{solution}{exo:g6:light-propagation:7}
À l'envers, et la gauche et la droite échangées~: sur l'écran, la
main s'agite en bas au lieu du haut, et du côté opposé. Les deux
retournements viennent du croisement des fils à la porte.
\end{solution}

\begin{solution}{exo:g6:light-propagation:8}
Les astronomes~: observer les éclipses sans danger sur l'écran plutôt
que dans le ciel. Les peintres~: décalquer la scène projetée pour
maîtriser la perspective. (Parents modernes~: tous les appareils
photo, et l'œil lui-même.)
\end{solution}

\begin{solution}{exo:g6:light-propagation:9}
L'\omterm{def:g1:light-and-shadows:shadow}{ombre} se resserre vers la taille réelle de la pièce~: la lampe
étant plus loin, les rayons frôlants la quittent moins inclinés ---
un éventail plus étroit --- et arrivent plus rapprochés. La
construction le montre~: des lignes plus plates, un écartement plus
serré.
\end{solution}

\begin{solution}{exo:g6:light-propagation:10}
Le Soleil est compact et lointain~: un seul éventail serré de rayons
quasi parallèles, une frontière bien nette --- des \omterm{def:g1:light-and-shadows:shadow}{ombres} franches.
Une couche de nuages gris est une lampe grande comme le ciel tout
entier~: la \omterm{def:g1:light-and-shadows:source}{lumière} arrive de toutes les directions à la fois, chaque
point du sol en reçoit une part, et les \omterm{def:g1:light-and-shadows:shadow}{ombres} s'effacent presque
complètement.
\end{solution}

\begin{solution}{exo:g6:light-propagation:11}
Plus lumineuse~: une porte large laisse entrer plus de fils venus de
chaque point de la scène. Plus floue~: à travers une porte large, les
fils de chaque point de la scène atteignent toute une \emph{tache}
d'écran au lieu d'un seul endroit --- d'innombrables taches qui se
chevauchent brouillent l'\omterm{prop:g5:mirrors-reflection:image}{image}. La netteté du sténopé \emph{est} sa
petitesse.
\end{solution}

\begin{solution}{exo:g6:light-propagation:12}
Chaque trouée entre les feuilles est un sténopé naturel~; chaque
tache lumineuse au sol est l'\omterm{prop:g5:mirrors-reflection:image}{image} du Soleil formée par ce
sténopé --- ronde un jour ordinaire, parce que le Soleil est rond.
Quand l'éclipse change le Soleil en croissant, chaque trouée projette
fidèlement le croissant~: des centaines de petites chambres noires,
un seul sujet céleste.
\end{solution}

\begin{solution}{pb:g6:light-propagation:1}
\textbf{1.} Contre l'écran, les rayons frôlants n'ont plus aucune
distance pour s'écarter~: l'\omterm{def:g1:light-and-shadows:shadow}{ombre} épouse le contour même de la
marionnette.
\textbf{2.} À mi-chemin, elle double~: \qty{40}{cm}.
\textbf{3.} Quatre fois $20$ font exactement \qty{80}{cm} --- cela
correspond~; la marionnette se tient au quart de la distance
lampe--écran.
\textbf{4.} Un large globe lumineux, ce sont plusieurs lampes à la
fois~: chaque partie jette sa propre \omterm{def:g1:light-and-shadows:shadow}{ombre} légèrement décalée, et les
bords du dragon s'étalent, doux et gris. Un théâtre net exige une
source petite, nue et quasi ponctuelle.
\textbf{5.} Vers la lampe~: l'éventail des rayons frôlants s'ouvre et
le loup grandit~; le bras reste bas, hors de l'éventail. À mesure que
la marionnette approche de la lampe, petite mais bien réelle,
l'adoucissement des bords augmente un peu --- de près, la largeur de
la lampe compte davantage.
\textbf{6.} Les lignes droites issues de la lampe et passant par les
deux marionnettes doivent aboutir sur des régions distinctes de
l'écran --- les marionnettes sont assez écartées latéralement pour
qu'aucune ne se trouve dans le cône d'\omterm{def:g1:light-and-shadows:shadow}{ombre} de l'autre.
\textbf{7.} Celle de la main~: plus près de la lampe, au tiers du
trajet, ses rayons frôlants s'écartent de trois fois sa taille, alors
que ceux du dragon, à mi-scène, ne doublent que --- près de la lampe,
même de petites mains jouent les géants.
\textbf{8.} Plaquer le méchant contre l'écran, et il se réduit à
lui-même (ou sort entièrement du \omterm{def:g6:light-propagation:ray}{faisceau})~; ou glisser un second
obstacle tout près de la lampe --- son \omterm{def:g1:light-and-shadows:shadow}{ombre} géante avale celle du
méchant, qui «~disparaît~» à l'intérieur d'une autre obscurité.
\textbf{9.} En bas du drap~: venu de la lampe haute, le fil
rectiligne passant par le trou continue vers le bas --- le sténopé
retourne.
\textbf{10.} De droite à gauche~: la porte échange la gauche et la
droite en même temps que le haut et le bas.
\textbf{11.} Plus lumineuse, mais désespérément floue~: chaque point
de la rue peint désormais une tache de la taille d'une pièce, les
taches se chevauchent et l'\omterm{prop:g5:mirrors-reflection:image}{image} se dissout --- le final meurt avec
la netteté.
\textbf{12.} «~Tout le programme de ce soir --- chaque \omterm{def:g1:light-and-shadows:shadow}{ombre}, chaque
géant, et une rue à l'envers --- a été joué par une seule loi~: la
\omterm{def:g1:light-and-shadows:source}{lumière} se propage en ligne droite.~»
\end{solution}
